\documentclass[12pt]{article}
\usepackage[utf8]{inputenc}
\usepackage[T1]{fontenc}
\usepackage[margin=1in]{geometry}
\usepackage{setspace}
\usepackage{xcolor}
\usepackage[colorlinks=true, urlcolor=blue!60!black, linkcolor=blue!60!black]{hyperref}
\setlength{\parindent}{0em}
\setlength{\parskip}{0.45em}
\renewcommand{\baselinestretch}{1.0}
\begin{document}
\begin{center}
{\Large \textit{Bivens v.\ Six Unknown Named Agents of Federal Bureau of Narcotics}, 403 U.S.\ 388 (1971)\par}
\vspace{0.3em}
{\normalsize Case Brief --- One-Page Edition\par}
\end{center}

\textbf{Memory Jogger.} Victims of unconstitutional federal searches may sue for damages under the Fourth Amendment itself (\href{https://supreme.justia.com/cases/federal/us/403/388/}{Justia}).

\textbf{Facts.} Federal narcotics agents entered Webster Bivens's apartment without a warrant or probable cause, manacled him before his family, searched the apartment ``from stem to stern,'' and strip-searched him at the courthouse.

\textbf{Procedural History.} With no statute authorizing suit against federal officers, the lower courts dismissed.

\textbf{Votes.} Reversed, 6--3: Brennan for the Court; Harlan concurred; Burger, Black, and Blackmun dissented.

\textbf{Holding.} A Fourth Amendment violation by federal agents supports a damages action absent special factors or a congressional substitute.

\textbf{Reasoning.} Federal power ``does not disappear like a magic gift when it is wrongfully used''; state tort law may be hostile to Fourth Amendment interests; damages are the traditional remedy. Harlan: for an innocent victim, ``it is damages or nothing.'' The dissenters left remedies to Congress. \emph{Later:} \textit{Abbasi} and \textit{Egbert} confined \textit{Bivens}.

\textbf{Examples.} \emph{Same-Side:} a warrantless DEA home raid. \emph{Opposite-Side:} a TSA checkpoint assault. \emph{Fence-Sitter:} an FBI fraud raid.

\textbf{Critique.} The case pits \textit{Marbury}'s maxim that every right needs a remedy against separation of powers, and its narrowing leaves federal agents far less accountable than state police under \S~1983.

\textbf{Key Quotation.} ``The present case involves no special factors counseling hesitation in the absence of affirmative action by Congress.'' --- Justice Brennan, 403 U.S.\ at 396.

\end{document}
